\documentclass[11pt]{article}
\usepackage[margin=1in]{geometry}
\usepackage{amsmath}
\usepackage{amssymb}
\usepackage{xcolor}
\usepackage{hyperref}
\usepackage{enumitem}
\usepackage{soul}
\usepackage{url}
\usepackage{booktabs}
\usepackage{graphicx}
\usepackage{caption}

% Margin relief: let LaTeX stretch interword spaces by up to 3em before
% emitting overfull-hbox warnings. Helps with long \texttt{} tokens
% (file paths, env var names) and prose-heavy paragraphs.
\usepackage{microtype}
\setlength{\emergencystretch}{8em}
\sloppy

% Reviewer comments shown in indented italics
\newenvironment{reviewer}{%
  \begin{quote}\itshape\color{black!70}}{%
  \end{quote}}

% Author responses shown plain
\newcommand{\response}[1]{\noindent\textbf{Response:} #1\par\medskip}

% Manuscript changes shown bulleted
\newenvironment{changes}{%
  \noindent\textbf{Changes to manuscript:}
  \begin{itemize}[leftmargin=2em,itemsep=2pt]}{%
  \end{itemize}\medskip}

\title{Response to Reviewers\\
\large Manuscript ID: ci-2026-009433.R1 (Second Revision)\\
``When Do Simple Models Win? Machine Learning Architectures for UV Absorption Prediction''}
\author{Umesh Arampath, Bryant Pero, David Stewart, David Demirjian}
\date{\today}

\begin{document}
\maketitle

We thank Associate Editor Prof.\ John Kitchin and both reviewers for their continued time and for the constructive second-round feedback, and we are grateful for the assessment that the work is ``nearly ready to be accepted.'' We respond below to the second-round comments: the editor's four technical issues and Reviewer~1's seven additional suggestions.

\clearpage

\section*{Response to Second-Round Comments}

\bigskip
\noindent\textbf{\large Editor's Technical Issues}
\medskip

\subsection*{Editor Comment E1: Table~5 caption recovery percentage}
\begin{reviewer}
The numbers in Table 5 caption don't appear to be consistent (the recovery \%) with the manuscript.
\end{reviewer}
\response{The reviewer is correct. Table~5's caption retained the ``93.6\% recovery'' figure we retracted in our first-round response to R1-15: the Mamede dataset is read directly from the authors' published Supporting Information rather than reconstructed by a Reaxys re-query, and RDKit canonicalises $74{,}783$ of $74{,}784$ entries ($\sim$100\%). The main text was corrected in the first revision; the caption was missed, and is now consistent.}
\begin{changes}
  \item \textbf{Table~5 caption.} Replaced ``our models use a reconstructed version (93.6\% recovery)'' with ``our models use a reconstruction from Mamede's published Supporting Information ($74{,}783$ of $74{,}784$ compounds canonicalised),'' consistent with the Photosafety Classification section and SI Section~S17.
\end{changes}

\subsection*{Editor Comment E2: Regenerate Figures 5, 6, and 9 with updated results}
\begin{reviewer}
it appears that figures 5, 6 and 9 need to be regenerated with the updated results
\end{reviewer}
\response{The reviewer is correct: Figures~5 (parity), 6 (error by wavelength and by solvent), and 9 (training time vs.\ performance) had been carried over from the v2 submission and not regenerated after the v3 (Greenman--Song) retrain. All three now come from the same v3 cross-validated predictions that underlie Table~2. One difference is worth flagging so it is not read as a discrepancy: Figures~5 and~6 \emph{pool} the test-set predictions across the five folds, giving a single pooled RMSE per model (RF~31.5, XGBoost~33.9, Chemprop~33.3, BiGRU~36.2, ChemBERTa~54.2~nm), whereas Table~2 and Figure~9 report the \emph{mean across folds}; the two agree to within $0.2$~nm for every model, and the Figure~5 caption now states which it shows. The ChemBERTa $\sim$556~nm parity ceiling (R1-16) and the physically-anchored axis limits (R2-10) are preserved, and the figure-generation scripts (\texttt{make\_parity\_plot.py}, \texttt{regenerate\_figures.py}) now read the v3 pooled predictions and are in the public repository.}
\begin{changes}
  \item \textbf{Figure~5 (parity plots).} Regenerated from v3 pooled predictions for all five model families; the per-panel RMSE/MAE/$R^2$ annotations are the pooled-prediction counterparts of the across-fold means reported in Table~2, agreeing with them to within $0.2$~nm.
  \item \textbf{Figure~6 (error breakdown).} Regenerated both panels (MAE by $\lambda_{\max}$ range; error distribution by solvent) from v3 predictions, aligned to the v3 dataset's solvent labels.
  \item \textbf{Figure~9 (training time vs.\ performance).} Regenerated with the v3 across-fold mean RMSE values reported in Table~2 (RF 31.5, XGBoost 33.9, Chemprop 33.1, BiGRU 36.2, ChemBERTa 54.1~nm).
\end{changes}

\subsection*{Editor Comment E3: Abstract vs.\ Methods on hyperparameter tuning}
\begin{reviewer}
The abstract suggests no hyperparameter tuning was used, but the methods suggest there was via a grid search
\end{reviewer}
\response{The reviewer is correct that this was contradictory. RF's hyperparameters (tree count and per-split feature fraction) were in fact selected by a 432-configuration grid search, so the abstract's and introduction's ``no hyperparameter tuning'' was inaccurate. We have reworded these statements to describe RF's tuning accurately while preserving the intended compute-cost contrast: unlike the deep models, RF needs only a lightweight one-time CPU grid search rather than GPU-based hyperparameter optimization.}
\begin{changes}
  \item \textbf{Abstract.} ``RF trains in 15 minutes on a CPU with no hyperparameter tuning'' $\rightarrow$ ``with only lightweight grid-search tuning (no GPU).''
  \item \textbf{Introduction (contributions bullet).} ``trains $\sim$100$\times$ faster with no hyperparameter tuning'' $\rightarrow$ ``trains $\sim$100$\times$ faster and needs only a lightweight CPU grid search rather than GPU-based hyperparameter optimization.''
  \item \textbf{Model Interpretability and Practical Considerations.} Reworded ``RF requires\ldots no hyperparameter tuning of learning rates or batch sizes'' to state that RF's only hyperparameters are set by a one-time CPU grid search, removing the implication that no tuning occurred.
\end{changes}

\subsection*{Editor Comment E4: Consistency of Figure~10 with Table~6}
\begin{reviewer}
Review if Figure 10 is consistent with Table 6.
\end{reviewer}
\response{The reviewer identified a genuine inconsistency. The submitted Figure~10 computed BiGRU saliency and its per-molecule prediction subtitles from the \emph{tuned} BiGRU ensemble, whereas Table~6 reports the \emph{default} BiGRU ensemble; the two disagreed (e.g., the old caption's Octocrylene ``+61~nm'' vs.\ Table~6's $+43$~nm). We regenerated Figure~10 from the same default-BiGRU 5-fold ensemble that Table~6 uses. The per-molecule predictions now match Table~6 exactly: Sinapic~Acid 318, Amiloxate 318, Padimate~O 330, Octisalate 333, Sulisobenzone 346, and Cinnamic~Acid 353~nm; the ensemble MAE (28.6~nm) matches Table~7. We updated the caption's example errors and the well/poorly-predicted thresholds to the regenerated values.}
\begin{changes}
  \item \textbf{Figure~10 (BiGRU saliency).} Regenerated from the default-BiGRU 5-fold ensemble; per-molecule Exp/Pred/error subtitles now match Table~6.
  \item \textbf{Figure~10 caption.} Updated the well/poorly-predicted thresholds ($<22 \rightarrow <20$~nm; $>28 \rightarrow >25$~nm) and the large-error example (``Octocrylene +61~nm, Cinnamic Acid +81~nm'' $\rightarrow$ ``Sulisobenzone +60~nm, Cinnamic Acid +78~nm'').
\end{changes}

\bigskip
\noindent\textbf{\large Reviewer 1 (Second Round)}
\medskip

\subsection*{Comment 1: Verbatim reproduction of reviewer comments}
\begin{reviewer}
The reviewer responses document says on page 1: ``each reviewer comment is reproduced verbatim in indented italics''. This is true for my comments R1-1 to R1-11, but R1-12 to R1-19 are paraphrased and sometimes include incorrect references to pages or line numbers. I think the paraphrasing still captures the meaning of my original comments, and in some cases is more detailed. So this is not inherently problematic, but this discrepancy leads me to suspect that this response document generation may have been AI-assisted. The authors should double check that the paraphrased reviewer comments are the only things that the AI didn't copy correctly.
\end{reviewer}
\response{We apologise for the discrepancy. The reviewer is correct that comments R1-12--R1-19 were paraphrased, with some incorrect page and line references. On the double-check the reviewer requested, we found that Reviewer~2's comments R2-1, R2-2, R2-4, R2-5, R2-6 and R2-7 had been paraphrased as well; every reviewer-comment block has now been re-verified against the original decision letter and the paraphrased ones corrected. Beyond the comment blocks we re-audited the response content and the manuscript, and the only substantive inaccuracies we found are the two factual corrections already reported in our first-round response: the ``variance is intrinsic to Chemprop'' claim, corrected via the R1-12 hyperparameter probe, and the ``93.6\% Reaxys recovery'' claim, corrected in R1-15. We found no other mis-transcribed content. On the reviewer's inference: yes, AI tools were used to help draft and format both the manuscript and the response document (see point~5); all analyses and responses are the authors' own and have been verified against the original decision letter and our data.}
\begin{changes}
  \item \textbf{Response document.} All reviewer-comment blocks re-checked against the original decision letter and the paraphrased ones (R1-12--R1-19; R2-1, R2-2, R2-4, R2-5, R2-6, R2-7) corrected. This document responds to the second-round comments; the corrected first-round responses can be supplied on request.
\end{changes}

\subsection*{Comment 2: Supporting Information section numbering}
\begin{reviewer}
References to the supporting information from the main text are of the form ``Supporting Information Section S7'', but the sections in the actual supporting information file are not numbered.
\end{reviewer}
\response{The reviewer is correct: the Supporting Information used achemso's default of \emph{unnumbered} sections, so the main-text ``Section~S$n$'' references pointed at sections that displayed no number. We have enabled numbering with an ``S'' prefix, so the SI now runs S1--S21, matching its already ``S''-prefixed tables and figures. We also checked every main-text cross-reference against the numbered SI and corrected an off-by-one that had gone unnoticed while the sections were unnumbered; all now resolve to the correct section.}
\begin{changes}
  \item \textbf{Supporting Information preamble.} Added \texttt{\textbackslash SectionNumbersOn} and \texttt{\textbackslash renewcommand\{\textbackslash thesection\}\{S\textbackslash arabic\{section\}\}} so sections render as S1--S21.
  \item \textbf{Main text.} Corrected every ``Supporting Information Section~S$n$'' cross-reference to match the numbered SI (e.g., the non-local-subset analysis is Section~S8, the ChemBERTa-ceiling diagnostic is Section~S21).
\end{changes}

\subsection*{Comment 3: Clearer definition of the v2/v3 dataset versions}
\begin{reviewer}
The new ``v2'' and ``v3'' descriptions of datasets used through the main text and SI could be more clearly defined, especially in their first usage. I also encourage the authors to consider using more descriptive titles for the dataset versions.
\end{reviewer}
\response{We have added an explicit definition at the point where the cleaned dataset is first produced (Methods, Dataset Description), naming the two versions descriptively: the \emph{v3} dataset is the ``Greenman--Song spread-filtered'' dataset ($18{,}415$ records) and the \emph{v2} dataset is the earlier ``first-occurrence-retention'' dataset ($18{,}755$ records), with a note that the two differ only in duplicate handling. We also added a parenthetical clarifying that these dataset-version labels are unrelated to the ``Chemprop v2'' software, since both terms appear in the manuscript.}
\begin{changes}
  \item \textbf{Methods (Dataset Description).} Added a sentence at first use defining v3 (Greenman--Song spread-filtered, $18{,}415$) and v2 (first-occurrence-retention, $18{,}755$), noting they differ only in duplicate handling, and disambiguating the labels from the Chemprop v2 software.
\end{changes}

\subsection*{Comment 4: Unconventional inter-revision phrasing in captions}
\begin{reviewer}
Some places in the manuscript mention how things were changed from one revision to the next, which is unconventional, e.g.:
- Table S16 caption: ``cleaned protocol adopted for the revised manuscript''
- Table S24 caption also appears to be missing a word: ``(adopted in this revision per )''
\end{reviewer}
\response{Corrected. We removed the revision-referencing phrasing from the two SI table captions the reviewer identified. The Table~S24 caption additionally contained a broken citation---the ``(adopted in this revision per )'' phrase ended with an empty \texttt{\textbackslash cite} command---which we have replaced with a proper citation to Greenman and Song. We also swept the manuscript and SI for other inter-revision phrasing and found no further instances.}
\begin{changes}
  \item \textbf{SI Table~S16 caption.} ``the Greenman--Song-style cleaned protocol adopted for the revised manuscript'' $\rightarrow$ ``the Greenman--Song-style spread-filtered protocol.''
  \item \textbf{SI Table~S24 caption.} ``The Greenman--Song cleaning rule (adopted in this revision per~) discards\ldots'' $\rightarrow$ ``The Greenman--Song cleaning rule~\cite{greenman2022multifidelity,song2025fluortools} discards\ldots'' (revision-mention removed, empty citation fixed).
\end{changes}

\subsection*{Comment 5: AI-assistance disclosure}
\begin{reviewer}
Acknowledgement / AI Disclosure: ``Claude Opus 4.6 (Anthropic) was used to assist with code debugging, code review, and manuscript formatting; all scientific content, analysis, and interpretation are solely the work of the authors.'' There is a commit on the GitHub repo titled ``De-AI prose cleanup: convert em-dashes to parens/commas; drop filler sentence-starters''. If AI assisted with any of the manuscript writing/revising in addition to formatting, I encourage the authors to add this to the disclosure. Similarly, the CLAUDE.md file seems to indicate that AI was used more heavily than just for code debugging/review. This is fine, but the disclosure statement should be updated if that's the case.
\end{reviewer}
\response{We appreciate the reviewer's attention to this, and the reviewer is correct that AI assistance extended beyond code debugging/review/formatting to editing of the manuscript text and this response document. We have updated the Acknowledgement to state this accurately. To be fully transparent, we have also left the repository intact---including the \texttt{CLAUDE.md} file and the commit history the reviewer references---rather than editing it. The scientific content, experimental work, data analysis, and interpretation are the authors' own.}
\begin{changes}
  \item \textbf{Acknowledgement.} Replaced ``\ldots used to assist with code debugging, code review, and manuscript formatting; all scientific content, analysis, and interpretation are solely the work of the authors'' with ``Claude (Anthropic) was used to assist with code development and debugging, data-analysis scripting, and editing of the manuscript and response-to-reviewer text; all scientific content, analysis, experimental work, and interpretation are solely the work of the authors, who take responsibility for the final manuscript.''
\end{changes}

\subsection*{Comment 6: Repository cleanup and LICENSE file}
\begin{reviewer}
I thank the authors for including the link to the public GitHub repo with the revision. The repo appears to be a comprehensive account of the project, even including many of the .tex files and figures, which is excellent. However, I suggest that the authors clean up the repository a bit by creating more subdirectories so all of the .tex files, images, etc. aren't in the project root. I think this will enhance readability and reuse of the codebase. Additionally, I noticed that while the authors specify in a couple places in the README that their project is published under the MIT License, there is no LICENSE file present in the repo, so this should be added to follow the standard practice.
\end{reviewer}
\response{Thank you for both suggestions. Rather than only rearranging our working repository, we have created a dedicated repository for the code accompanying this manuscript, \url{https://github.com/Umesh1608/ML_UV_models_jcim}, so that the work has a single, cleanly organised point of reference. It is MIT licensed with a \texttt{LICENSE} file at its root, and it contains only what is needed to understand and reproduce the study: the \texttt{uvml} model library; the runnable scripts grouped by purpose (\texttt{training}, \texttt{tuning}, \texttt{evaluation}, \texttt{analysis}, \texttt{data\_prep}, \texttt{figures}); the processed datasets; and the per-fold cross-validation metrics. Nothing sits in the project root except the README, the \texttt{LICENSE}, and the packaging files. The Data and Software Availability section now cites this repository.

Assembling it also let us find and correct two reproducibility defects that our earlier root-level-to-subdirectory reorganisation had introduced, which we report here for completeness: the scripts resolved \texttt{data/} and \texttt{results/} relative to their own subdirectory rather than the repository root, so they raised \texttt{FileNotFoundError} on a fresh clone; and the README documented command-line flags that the scripts do not accept. Both are fixed. The per-fold cross-validation metrics are now committed to the repository, so the RMSE values in Table~2 can be reproduced from a clone with a single command, without any download, GPU, or retraining.

Our original development repository (\url{https://github.com/Umesh1608/ML_UV_models}) remains public and unchanged, including the \texttt{CLAUDE.md} file and the commit history the reviewer refers to in point~5 above; we have not edited or removed anything from it.}
\begin{changes}
  \item \textbf{New code repository.} Created \url{https://github.com/Umesh1608/ML_UV_models_jcim}: MIT licensed with a \texttt{LICENSE} file at the root, scripts grouped by purpose, processed datasets and per-fold metrics included, and no loose files in the project root.
  \item \textbf{Data and Software Availability (pg.\ 41).} Updated to cite the new repository and to note that the per-fold metrics it contains reproduce the Table~2 values directly from a clone.
  \item \textbf{Reproducibility fixes in the released code.} Scripts now resolve \texttt{data/} and \texttt{results/} against the repository root, and the README documents the flags the scripts actually accept.
  \item \textbf{Original repository.} Left public and unchanged, including \texttt{CLAUDE.md} and the full commit history.
\end{changes}

\subsection*{Comment 7: Figure~2(b) parallel vs.\ sequential D-MPNN encoders}
\begin{reviewer}
Figure 2(b): The Chemprop D-MPNN part of this figure is misleading because it implies the Solute and Solvent D-MPNNs are sequential rather than parallel. This is explained correctly in the caption, but I think the figure should be updated accordingly.
\end{reviewer}
\response{Corrected. We redrew Figure~2(b) so the solute and solvent D-MPNN encoders are shown as two \emph{parallel} branches---both fed from the shared ``RDKit~$\to$~graphs'' step and both feeding the ``Mean-pool~$\oplus$~concat'' step---rather than as a vertical sequence in which one appears to feed the other. The figure now matches the multicomponent architecture described in the caption.}
\begin{changes}
  \item \textbf{Figure~2(b).} Replaced the sequential Solute-D-MPNN~$\rightarrow$~Solvent-D-MPNN layout with two side-by-side parallel encoder branches that merge at the mean-pool/concatenation step.
\end{changes}

\clearpage

We hope these revisions address the remaining concerns. We thank the Associate Editor and both reviewers for the rigour and constructiveness of their feedback across both rounds of review.

\bibliographystyle{plain}
\bibliography{Proposal}

\end{document}
